\documentclass[10pt]{beamer}

% --- Modern Theme ---
\usetheme{metropolis}
\usecolortheme{default}
\setbeamercolor{background canvas}{bg=white}

% --- Packages ---
\usepackage[utf8]{inputenc}
\usepackage{booktabs}
\usepackage{array}
\usepackage{tabularx}
\usepackage{amsmath}
\usepackage{amssymb}
\usepackage[scale=2]{ccicons}
\usepackage{pgfplots}
\usepgfplotslibrary{dateplot}
\usepackage{xspace}

% --- Title Information ---
\title{Foundations of Artificial Intelligence}
\subtitle{Part VIII: AI in Society, Ethics \& Emerging Trends}
\institute{Department of Computer Science}
\date{}

\begin{document}

\maketitle

\begin{frame}{Lecture Roadmap}
    \setbeamertemplate{section in toc}[sections numbered]
    \tableofcontents[hideallsubsections]
\end{frame}

\begin{frame}{Learning Objectives}
    By the end of this lecture, you should be able to:
    \begin{itemize}
        \item Identify and explain sources of bias in AI systems.
        \item Discuss privacy implications of data-driven AI.
        \item Analyse AI's impact on labour and employment.
        \item Articulate principles of Responsible AI and governance frameworks.
        \item Describe Generative AI and Large Language Models at a conceptual level.
        \item Recognise opportunities for AI to create social good, including agriculture and public policy.
    \end{itemize}
\end{frame}

\section{Bias in AI Systems}

\begin{frame}{Where Bias Comes From}
    Bias in AI arises at multiple stages:
    \begin{itemize}
        \item \textbf{Data collection:} datasets under-represent certain groups.
        \item \textbf{Label quality:} annotators bring implicit assumptions.
        \item \textbf{Model choice:} some architectures amplify existing patterns.
        \item \textbf{Deployment context:} model used outside the population it was trained on.
    \end{itemize}

    \vspace{0.2cm}
    \textbf{Core insight (from R\&N, Ch.~26):}
    An AI system optimises for measurable proxies; if those proxies are biased, the system inherits---and often amplifies---that bias.
\end{frame}

\begin{frame}{Types of Bias}
    \small
    \setlength{\tabcolsep}{7pt}
    \renewcommand{\arraystretch}{1.25}
    \begin{tabularx}{\textwidth}{l X}
        \toprule
        \textbf{Type} & \textbf{Description} \\
        \midrule
        Historical bias    & Patterns in data reflect past discrimination. \\
        Representation bias & Certain groups appear less in training data. \\
        Measurement bias   & Features poorly measured for some sub-groups. \\
        Aggregation bias   & One model applied to diverse populations. \\
        Deployment bias    & System used in context different from design intent. \\
        \bottomrule
    \end{tabularx}

    \vspace{0.2cm}
    Example: Facial recognition systems with higher error rates on darker skin tones (Buolamwini \& Gebru, 2018).
\end{frame}

\begin{frame}{Measuring Fairness}
    Multiple mathematical definitions of fairness:
    \begin{itemize}
        \item \textbf{Demographic parity:} equal positive prediction rates across groups.
        \item \textbf{Equal opportunity:} equal true positive rates across groups.
        \item \textbf{Calibration:} predicted probabilities match actual outcomes equally well.
    \end{itemize}

    \vspace{0.2cm}
    \textbf{Key result:} It is mathematically \emph{impossible} to satisfy all definitions simultaneously when base rates differ across groups (Chouldechova, 2017).

    \vspace{0.2cm}
    Choosing a fairness criterion is therefore an \textbf{ethical and policy decision}, not just a technical one.
\end{frame}

\begin{frame}{Mitigating Bias}
    Three intervention points:
    \begin{enumerate}
        \item \textbf{Pre-processing:} re-sample, re-weight, or augment training data.
        \item \textbf{In-processing:} add fairness constraints to the learning objective.
        \item \textbf{Post-processing:} adjust decision thresholds per group after training.
    \end{enumerate}

    \vspace{0.2cm}
    Complementary practices:
    \begin{itemize}
        \item Diverse teams building diverse datasets.
        \item Regular bias audits throughout the model lifecycle.
        \item Participatory design involving affected communities.
    \end{itemize}
\end{frame}

\section{Privacy Concerns}

\begin{frame}{AI and Personal Data}
    Modern AI is data hungry:
    \begin{itemize}
        \item Training on web-scale corpora may include personal information.
        \item Inference can reveal sensitive attributes not explicitly shared.
        \item Models can memorise and reproduce training samples.
    \end{itemize}

    \vspace{0.2cm}
    \textbf{From the textbook (Ch.~26):} ``Computerization will lead to loss of privacy.''

    \vspace{0.2cm}
    Privacy risk categories:
    \begin{itemize}
        \item Direct identification from outputs.
        \item Re-identification by combining predictions with external data.
        \item Membership inference attacks.
    \end{itemize}
\end{frame}

\begin{frame}{Privacy-Preserving Techniques}
    \begin{itemize}
        \item \textbf{Differential privacy:} add calibrated noise so individual records cannot be inferred.
        \item \textbf{Federated learning:} train on distributed data; model updates (not raw data) are shared.
        \item \textbf{Data minimisation:} collect and retain only what is strictly necessary.
        \item \textbf{Anonymisation:} remove identifiers before training (but beware re-identification).
    \end{itemize}

    \vspace{0.2cm}
    Regulatory context: GDPR (EU), CCPA (California), and emerging national frameworks set binding requirements on data use and consent.
\end{frame}

\section{AI and Employment}

\begin{frame}{The Automation Question}
    AI and automation historically:
    \begin{itemize}
        \item Displace \textbf{routine, repetitive} tasks (both cognitive and physical).
        \item Create new roles requiring human judgment, creativity, and adaptability.
    \end{itemize}

    \vspace{0.2cm}
    Recent evidence:
    \begin{itemize}
        \item Tasks structured around rules are most vulnerable.
        \item Jobs combining physical dexterity + social intelligence are harder to automate.
        \item Generative AI now extends reach to knowledge work and content creation.
    \end{itemize}

    \vspace{0.2cm}
    \textbf{R\&N perspective:} AI creates wealth but its distribution is not guaranteed to be equitable.
\end{frame}

\begin{frame}{AI's Impact on the Labour Market}
    \small
    \setlength{\tabcolsep}{7pt}
    \renewcommand{\arraystretch}{1.25}
    \begin{tabularx}{\textwidth}{l X X}
        \toprule
        \textbf{Effect} & \textbf{Description} & \textbf{Example} \\
        \midrule
        Displacement & Tasks automated, jobs eliminated & Data entry, basic translation \\
        Augmentation & Human + AI outperforms either alone & Radiology + AI screening \\
        Creation & New roles emerge around AI systems & Prompt engineering, AI auditing \\
        Polarisation & Demand shifts to high- and low-skill extremes & Middle-skill roles most at risk \\
        \bottomrule
    \end{tabularx}

    \vspace{0.2cm}
    Outcome depends on transition support, education policy, and regulation.
\end{frame}

\begin{frame}{Adapting to AI-Driven Change}
    Strategies for society and individuals:
    \begin{itemize}
        \item \textbf{Continuous upskilling:} lifelong learning and digital literacy.
        \item \textbf{Education reform:} prioritise critical thinking, creativity, and interdisciplinary skills.
        \item \textbf{Policy instruments:} social safety nets, portable benefits, retraining programmes.
        \item \textbf{Complementarity:} design AI tools that amplify human capabilities.
    \end{itemize}
\end{frame}

\section{Responsible AI}

\begin{frame}{What Is Responsible AI?}
    Responsible AI is a set of \textbf{principles and practices} to ensure AI systems are:

    \vspace{0.15cm}
    \begin{itemize}
        \item \textbf{Fair:} do not discriminate unjustly.
        \item \textbf{Reliable \& Safe:} behave correctly under normal and edge-case conditions.
        \item \textbf{Transparent:} decisions are explainable to users and regulators.
        \item \textbf{Privacy-preserving:} handle personal data with care.
        \item \textbf{Inclusive:} accessible to and beneficial for diverse populations.
        \item \textbf{Accountable:} humans can identify and correct failures.
    \end{itemize}
\end{frame}

\begin{frame}{Explainability and Interpretability}
    \textbf{Interpretable model:} inherently understandable (e.g., decision tree, logistic regression).

    \textbf{Explainable model:} post-hoc explanations for black-box models.

    \vspace{0.2cm}
    Explanation methods:
    \begin{itemize}
        \item \textbf{LIME:} locally approximate model behaviour around a specific prediction.
        \item \textbf{SHAP:} attribute prediction contribution to each input feature.
        \item \textbf{Attention maps:} visual evidence for image/text model decisions.
        \item \textbf{Counterfactuals:} ``If feature X had been different, the outcome would have changed.''
    \end{itemize}

    \vspace{0.2cm}
    Explainability is not optional in \textbf{high-stakes domains}: healthcare, criminal justice, credit scoring.
\end{frame}

\begin{frame}{Responsibility and Accountability}
    A challenge highlighted in foundational AI literature:

    \begin{quote}
        \small ``One particular risk is \textbf{loss of accountability}: if an AI system makes a bad decision, who is responsible?''
    \end{quote}

    \vspace{0.1cm}
    Stakeholders involved:
    \begin{itemize}
        \item Data providers --- responsible for data quality and consent.
        \item Developers --- responsible for model design and testing.
        \item Deployers --- responsible for context and operational safeguards.
        \item Regulators --- responsible for setting standards and enforcement.
        \item Users --- responsible for appropriate use.
    \end{itemize}
\end{frame}

\section{AI Governance}

\begin{frame}{Why AI Governance?}
    Without governance, AI systems can:
    \begin{itemize}
        \item Perpetuate or amplify social inequalities.
        \item Operate in domains requiring oversight without adequate checks.
        \item Pose safety risks in high-stakes environments.
        \item Concentrate economic and political power.
    \end{itemize}

    \vspace{0.2cm}
    Governance brings:
    \begin{itemize}
        \item Rules and standards that constrain harmful deployment.
        \item Incentives for responsible development.
        \item Redress mechanisms for affected individuals.
    \end{itemize}
\end{frame}

\begin{frame}{Governance Frameworks (Global Overview)}
    \small
    \setlength{\tabcolsep}{6pt}
    \renewcommand{\arraystretch}{1.25}
    \begin{tabularx}{\textwidth}{l X}
        \toprule
        \textbf{Framework / Region} & \textbf{Key Approach} \\
        \midrule
        EU AI Act (2024)     & Risk-based tiering; high-risk systems face strict requirements; bans some uses. \\
        OECD AI Principles   & Human-centred values, transparency, robustness, accountability, inclusiveness. \\
        US Executive Order (2023) & Safety testing, watermarking generative AI, federal agency guidance. \\
        UNESCO AI Ethics     & Global recommendation on data governance, environment, education, and gender. \\
        National frameworks  & Countries developing their own policies; divergence in approach is a challenge. \\
        \bottomrule
    \end{tabularx}
\end{frame}

\begin{frame}{Technical Governance Tools}
    \begin{itemize}
        \item \textbf{Model cards:} document model purpose, performance across sub-groups, and intended use.
        \item \textbf{Datasheets for datasets:} record provenance, collection methods, and known limitations.
        \item \textbf{Algorithmic impact assessments:} structured process to evaluate potential harms before deployment.
        \item \textbf{Audit trails and logging:} maintain records of model decisions for later review.
        \item \textbf{Red-teaming:} adversarial testing by independent teams to find safety issues.
    \end{itemize}
\end{frame}

\section{Generative AI \& Large Language Models}

\begin{frame}{What Is Generative AI?}
    Generative AI produces \textbf{new content} (text, images, code, audio, video) rather than only classifying existing content.

    \vspace{0.2cm}
    Key model families:
    \begin{itemize}
        \item \textbf{Autoregressive models (GPT):} predict next token given context.
        \item \textbf{Diffusion models:} learn to reverse a noise process to generate images.
        \item \textbf{VAEs and GANs:} encode and decode latent representations.
    \end{itemize}

    \vspace{0.2cm}
    Common applications: writing assistance, code generation, image synthesis, music, and video.
\end{frame}

\begin{frame}{Large Language Models (LLMs)}
    \textbf{LLMs} are autoregressive transformers trained on internet-scale text.

    \vspace{0.2cm}
    \textbf{Key capability:} Predict
    \[
        P(\text{token}_t \mid \text{token}_{1:t-1})
    \]
    across billions of parameters, learning grammar, facts, and reasoning patterns.

    \vspace{0.2cm}
    \textbf{Emergent abilities at scale:}
    \begin{itemize}
        \item Few-shot and zero-shot task completion.
        \item Step-by-step reasoning (chain-of-thought prompting).
        \item Code generation, translation, and summarisation.
    \end{itemize}
\end{frame}

\begin{frame}{How LLMs Work (Conceptual)}
    Three-stage development cycle:
    \begin{enumerate}
        \item \textbf{Pre-training:} self-supervised next-token prediction on large corpora.
        \item \textbf{Fine-tuning (SFT):} supervised training on curated instruction-response pairs.
        \item \textbf{Alignment (RLHF):} human feedback used to improve helpfulness and safety.
    \end{enumerate}

    \vspace{0.2cm}
    \textbf{Transformer architecture:}
    \begin{itemize}
        \item Self-attention: each token attends to all others in context.
        \item Feed-forward layers: apply non-linear transformations.
        \item Scales well with data and compute.
    \end{itemize}
\end{frame}

\begin{frame}{Capabilities and Limitations of LLMs}
    \textbf{Capabilities:}
    \begin{itemize}
        \item Broad knowledge across domains from training data.
        \item Flexible generation: adapt tone, style, and format on request.
        \item Strong in-context learning without retraining.
    \end{itemize}

    \vspace{0.2cm}
    \textbf{Limitations:}
    \begin{itemize}
        \item \textbf{Hallucination:} confidently generate false information.
        \item \textbf{No real-time knowledge:} knowledge cut-off from training.
        \item \textbf{Bias amplification:} reproduce and magnify biases in training data.
        \item \textbf{Misuse risk:} disinformation, phishing, manipulation at scale.
        \item \textbf{Computational cost:} training and inference are energy-intensive.
    \end{itemize}
\end{frame}

\section{AI for Social Good}

\begin{frame}{AI for Social Good}
    AI creates opportunities to address large-scale societal challenges:
    \begin{itemize}
        \item \textbf{Healthcare:} disease detection, drug discovery, diagnostic support.
        \item \textbf{Education:} personalised learning, accessibility tools, tutoring.
        \item \textbf{Environment:} climate modelling, wildlife conservation, energy optimisation.
        \item \textbf{Disaster response:} satellite imagery analysis, resource coordination.
        \item \textbf{Legal/civic:} document summarisation, access to information.
    \end{itemize}

    \vspace{0.15cm}
    Success depends on: \textbf{quality data}, \textbf{community involvement}, and \textbf{sustainable deployment}.
\end{frame}

\begin{frame}{AI in Agriculture}
    Agriculture accounts for a significant share of employment and food security in many countries.

    \vspace{0.2cm}
    How AI can help:
    \begin{itemize}
        \item \textbf{Crop disease detection:} image classification on leaf photos to diagnose disease early.
        \item \textbf{Yield prediction:} regression models using weather, soil, and satellite data.
        \item \textbf{Precision irrigation:} sensor data + ML to water only where and when needed.
        \item \textbf{Pest and weed management:} targeted treatment reduces chemical use.
        \item \textbf{Market price prediction:} helps smallholder farmers plan and reduce losses.
    \end{itemize}

    \vspace{0.15cm}
    \textit{Mobile-first, low-bandwidth models are critical for low-resource settings.}
\end{frame}

\begin{frame}{AI in Public Policy}
    Governments are adopting AI across policy domains:
    \begin{itemize}
        \item \textbf{Public health:} epidemiological modelling, outbreak detection, resource allocation.
        \item \textbf{Transport:} traffic prediction, public transit optimisation.
        \item \textbf{Social services:} fraud detection in benefit systems, eligibility triage.
        \item \textbf{Urban planning:} predictive analysis for infrastructure and housing needs.
        \item \textbf{Tax and revenue:} anomaly detection in filings.
    \end{itemize}

    \vspace{0.2cm}
    \textbf{Risks:} automated decisions affecting citizens require transparency, appeal mechanisms, and human oversight.
\end{frame}

\begin{frame}{Power and Responsibility}
    AI systems concentrates significant power:
    \begin{itemize}
        \item Economic: platform lock-in, IP concentration.
        \item Political: surveillance capability, information shaping.
        \item Social: algorithmic gatekeeping of opportunity.
    \end{itemize}

    \vspace{0.2cm}
    \textbf{Countervailing forces:}
    \begin{itemize}
        \item Open-source models reduce barriers to access.
        \item Civil society audits and advocacy.
        \item International coordination on standards.
        \item Education and public AI literacy.
    \end{itemize}

    \vspace{0.15cm}
    \textbf{Core principle:} The greater the capability, the greater the obligation to deploy responsibly.
\end{frame}

\begin{frame}{Live Demo (In-Class Activity)}
    \textbf{Interactive file:} \texttt{ai-ethics-governance-demo.html}

    \vspace{0.15cm}
    Suggested walkthrough (10 min):
    \begin{enumerate}
        \item \textbf{Bias and Fairness}: adjust group outcomes and interpret parity gap.
        \item \textbf{Privacy Risk}: vary noise and retention to show utility-risk trade-off.
        \item \textbf{Employment Impact}: compare displacement pressure vs augmentation.
        \item \textbf{Governance Check}: run checklist to decide deployment readiness.
        \item \textbf{Generative AI Risk}: scan prompts for policy and misuse signals.
    \end{enumerate}

    \vspace{0.15cm}
    \small
    \textbf{URL:} \url{http://127.0.0.1:8000/Lecture%208/ai-ethics-governance-demo.html}

    \vspace{0.1cm}
    \textit{Run from workspace root:} \texttt{python3 -m http.server 8000}
\end{frame}

\begin{frame}{Summary and Course Conclusion}
    \textbf{What we covered today:}
    \begin{itemize}
        \item Bias in AI systems and fairness trade-offs.
        \item Privacy risks and mitigation techniques.
        \item AI's nuanced effects on employment.
        \item Responsible AI principles and explainability.
        \item AI governance frameworks and tools.
        \item Generative AI and LLMs: power and pitfalls.
        \item Social good applications: agriculture and public policy.
    \end{itemize}

    \vspace{0.2cm}
    \textbf{Key takeaway:}
    Technical capability must always be paired with ethical reflection, inclusive design, and accountable governance.

    \vspace{0.15cm}
    \textit{``The question is not whether we will have AI---it is what kind of AI we choose to build, and for whom.''}
\end{frame}

\end{document}
